\documentclass[11pt]{article}

\usepackage[margin=1.15in]{geometry}
\usepackage{amsmath,amssymb,amsthm}
\usepackage{microtype}

\theoremstyle{plain}
\newtheorem{theorem}{Theorem}[section]
\newtheorem{proposition}[theorem]{Proposition}
\newtheorem{lemma}[theorem]{Lemma}
\newtheorem{corollary}[theorem]{Corollary}
\theoremstyle{definition}
\newtheorem{definition}[theorem]{Definition}
\theoremstyle{remark}
\newtheorem{remark}[theorem]{Remark}

\newcommand{\R}{\mathbb{R}}
\newcommand{\Q}{\mathbb{Q}}
\newcommand{\E}{\mathbb{E}}
\newcommand{\grad}{\nabla}

\title{The loss landscape of a single polynomial neuron}
\author{Andrew Mullhaupt}
\date{\today}

\begin{document}
\maketitle

\begin{abstract}
For the least-squares loss of a one-neuron polynomial network the phase
portrait of the gradient flow can be determined exactly: the critical set is
algebraic, the Morse property is decidable by polynomial gcd over $\Q$, and
the merge tree of the sublevel sets follows from a reduction of the loss to a
one-dimensional potential along a \emph{backbone} curve. What does not follow
is the connectivity of the invariant manifolds, because a saddle connection is
codimension one. We exhibit a one-parameter family --- a similarity of the
plane --- along which every algebraic invariant of the landscape is constant
while the connectivity nevertheless bifurcates, which shows that no elimination
argument can decide the question and places it instead in the Melnikov setting.
We then describe what can be certified numerically, and what the certified
picture says about the descent methods actually used to train such a model.
\end{abstract}

\section{The model}

Let $x$ be a real random variable whose raw moments
\(
\mu_k = \E[x^k]
\)
exist for $k = 0,\dots,2D$, and let $f$ and $g$ be real polynomials of degrees
$d_f$ and $d_g$ with $D = \max(d_f, d_g)$. The network is the single neuron
\[
  F(x; a, b) \;=\; a\, g(bx),
\]
with an outer weight $a$ and an inner weight $b$, and the objective is the
population least-squares loss
\begin{equation}\label{eq:loss}
  L(a,b) \;=\; \E\bigl[(f(x) - a\,g(bx))^2\bigr]
         \;=\; C - 2a\,B(b) + a^2 A(b),
\end{equation}
where
\begin{equation}\label{eq:ABC}
  A(b) = \E\bigl[g(bx)^2\bigr], \qquad
  B(b) = \E\bigl[f(x)\,g(bx)\bigr], \qquad
  C    = \E\bigl[f(x)^2\bigr].
\end{equation}
Writing $f = \sum_i f_i x^i$ and $g = \sum_j g_j x^j$ and taking expectations
term by term,
\begin{equation}\label{eq:ABCcoeffs}
  A(b) = \sum_{i,j} g_i g_j \mu_{i+j}\, b^{\,i+j}, \qquad
  B(b) = \sum_{j} g_j \Bigl( \sum_i f_i \mu_{i+j} \Bigr) b^{\,j}, \qquad
  C    = \sum_{i,k} f_i f_k \mu_{i+k},
\end{equation}
so $A$ has degree $2d_g$, $B$ has degree $d_g$, and the problem depends on the
input law \emph{only} through the finite moment vector
$(\mu_0,\dots,\mu_{2D})$. Two consequences are worth stating at once. First,
the landscape is a property of the moments, not of the distribution: any two
laws agreeing to order $2D$ give the same $L$. Second, if the moments are
rational --- as they are for $x \sim U(0,1)$, where $\mu_k = 1/(k+1)$, and for
any empirical law on rational samples --- then $A$, $B$ and $C$ lie in
$\Q[b]$ and every statement below about the critical set is decidable by exact
arithmetic.

Throughout we assume

\begin{quote}
\textbf{(A)} \quad $A(b) \neq 0$ for every real $b$.
\end{quote}

Since $A$ is a polynomial, (A) says exactly that $A$ has constant sign; after
replacing $\mu$ by $-\mu$ if necessary we may and do assume $A > 0$. When $\mu$
is the moment sequence of a positive measure and $g \not\equiv 0$, (A) holds
automatically, because $A(b) = \E[g(bx)^2]$. The hypothesis has content only
when the moment functional is indefinite; see Remark~\ref{rem:krein}.

\begin{remark}\label{rem:krein}
A finite real sequence $(\mu_k)$ need not be the moment sequence of a positive
measure. It is always the moment sequence of a \emph{signed} measure, and
signed quadrature is ordinary numerical practice --- closed Newton--Cotes has
negative weights from order nine. In that case $\langle p, q\rangle =
\E[p(x)q(x)]$ is an indefinite bilinear form, a Krein rather than a Hilbert
structure, and $L$ can take negative values. Everything below survives, because
nothing below uses $L \geq 0$: it uses only (A), which asks the form to be
positive on the one-dimensional span of $g(b\,\cdot)$ for each real $b$. That
can hold while the ambient Hankel matrix of $\mu$ is indefinite, the negative
cone simply missing the curve $b \mapsto (g_0, g_1 b, \dots, g_{d_g}b^{d_g})$.
What is lost in that case is not the geometry but the interpretation of $L$ as
an approximation error.
\end{remark}

\section{The backbone reduction}

For fixed $b$, \eqref{eq:loss} is a quadratic in $a$ with leading coefficient
$A(b) > 0$, so it has a unique minimiser. Completing the square gives the
identity on which everything else rests.

\begin{proposition}[Backbone form]\label{prop:backbone}
Set
\[
  a^*(b) = \frac{B(b)}{A(b)}, \qquad
  u(b)   = C - \frac{B(b)^2}{A(b)}, \qquad
  w      = a - a^*(b).
\]
Then $L(a,b) = u(b) + A(b)\,w^2$.
\end{proposition}

We call the curve $\{w = 0\}$, i.e.\ $a = a^*(b)$, the \emph{backbone}, and
$u$ the \emph{backbone potential}. The level sets of $L$ are symmetric about
the backbone in the $a$ direction, a fact we use in \S\ref{sec:skeleton}.

\begin{proposition}[Critical structure]\label{prop:critical}
Assume \textup{(A)}. Then:
\begin{enumerate}
\item[(i)] $L_{aa} = 2A > 0$ everywhere, so the Hessian of $L$ is nowhere
  negative semidefinite: $L$ has no local maximum in $\R^2$.
\item[(ii)] Every critical point lies on the backbone.
\item[(iii)] On the backbone, $\det \operatorname{Hess} L = 2A\,u''$; hence a
  critical point is a minimum when $u'' > 0$, a saddle when $u'' < 0$, and
  degenerate exactly when $u'' = 0$.
\item[(iv)] With $N := A'B - 2B'A$,
  \begin{equation}\label{eq:uprime}
    u' \;=\; \frac{B\,N}{A^2},
  \end{equation}
  so the critical values of $b$ are precisely the real roots of $B\cdot N$.
\end{enumerate}
\end{proposition}

\begin{proof}
$L_a = 2(aA - B) = 2Aw$, which vanishes iff $w = 0$; this is (i) and (ii),
the first because $L_{aa} = 2A > 0$ independently of the other entries.
For (iii), differentiate $L = u + Aw^2$ with $w = a - a^*(b)$. At $w = 0$ one
finds $L_{aa} = 2A$, $L_{ab} = -2A\,a^{*\prime}$ and $L_{bb} = u'' + 2A\,
a^{*\prime 2}$, whence
\[
  \det\operatorname{Hess}L = 2A\bigl(u'' + 2A a^{*\prime 2}\bigr) - 4A^2
  a^{*\prime 2} = 2A\,u''.
\]
For (iv), differentiate $u = C - B^2/A$:
\[
  u' = -\frac{2BB'A - B^2A'}{A^2} = \frac{B\,(A'B - 2B'A)}{A^2}
     = \frac{B\,N}{A^2}. \qedhere
\]
\end{proof}

Two features of \eqref{eq:uprime} deserve emphasis, since both are used later.

\paragraph{The roof.} At a real root $b_0$ of $B$ we have $a^*(b_0) = 0$ and
$u(b_0) = C$. \emph{All} roots of $B$ therefore carry the same critical value
$C = \E[f^2]$. Moreover, if $b_0$ is a simple root then
$N(b_0) = -2B'(b_0)A(b_0)$ and $u''(b_0) = -2B'(b_0)^2/A(b_0) < 0$, so
\begin{equation}\label{eq:Bsaddle}
  \det \operatorname{Hess} L\big|_{b_0} \;=\; -4B'(b_0)^2 ,
\end{equation}
and every simple root of $B$ is a saddle, at the common height $C$. A strictly
descending trajectory therefore cannot run from one root of $B$ to another.

\paragraph{Alternation.} By (iii) the type of a critical point is the sign of
$u''$ at a critical point of the one-dimensional function $u$. If $u$ is Morse
its critical points alternate between local minima and local maxima along the
$b$-axis, so the critical points of $L$ alternate between minima and saddles.
The two-dimensional landscape inherits the combinatorics of a
one-dimensional one.

\section{Deciding the Morse property exactly}

Let $H := B \cdot N$, so that by \eqref{eq:uprime} the critical $b$-values are
the real roots of $H$ and, $A$ having no real root, $u'$ vanishes to exactly
the order to which $H$ does.

\begin{proposition}\label{prop:morse}
Assume \textup{(A)} and $H \not\equiv 0$. Then $L$ is Morse if and only if $H$
has no repeated real root, equivalently if and only if $\gcd(H, H')$ has no
real root.
\end{proposition}

\begin{proof}
A critical point is degenerate iff $u'' = 0$ there, by
Proposition~\ref{prop:critical}(iii). If $b_0$ is a root of $H$ of multiplicity
$m$ then $u' = H/A^2$ vanishes to order $m$ at $b_0$, so $u''(b_0) = 0$ iff
$m \geq 2$.
\end{proof}

Both tests are effective over $\Q$, and it is worth saying how, since the
usual expectation would be a numerical root-finder followed by a tolerance.
There is none here: no root of $H$ is ever computed as a floating-point
number for the purpose of deciding anything.

Clear denominators so that $H \in \mathbb{Z}[b]$. The Sturm chain is the
Euclidean remainder sequence $H_0 = H$, $H_1 = H'$,
$H_{k+1} = -\operatorname{rem}(H_{k-1}, H_k)$, each remainder replaced by its
primitive part --- the integer polynomial obtained by dividing out the content
--- which is a positive rational rescaling and therefore changes no sign the
theory reads. Removing the content at every step is what keeps the chain in
manageable integers. The textbook alternative, the subresultant PRS, avoids
computing the content at the price of carrying larger coefficients; on these
polynomials that price is measured and substantial, the peak intermediate
coefficient reaching 20116 bits against 8257 for $A^2$ and 39554 against 8573
for $A^4$. Primitive PRS is the coefficient-optimal variant here. The two
chains have identical sign sequences --- dividing out the content multiplies
by a positive rational --- so the choice is one of space and time and not of
correctness.

Sturm's theorem then gives the number of distinct real roots of $H$ in any
interval $(p, q]$ with rational endpoints as the drop in sign variations of
the chain evaluated at $p$ and $q$. A Cauchy bound on the coefficients
supplies a starting interval containing every real root; repeated bisection,
using the same count, isolates each root in an interval containing exactly one
and refines it to any rational width. The output is therefore a list of
rational intervals, each certified to hold exactly one real root of $H$ --- not
a list of approximate roots.

Every subsequent decision is an exact sign evaluation on those intervals. The
type of a critical point is the sign of $u''$ there, obtained from the sign of
a polynomial on the isolating interval rather than from the value of $u''$ at
an approximate root. The ordering of critical values, which determines the
merge tree, is decided the same way: a candidate rational level $c$ separates
two critical values when the level polynomial $(C-c)A - B^2$ has the
appropriate root count on each side, so critical values are never compared to
one another numerically and coincident values --- which are generic here,
since every root of $B$ has value exactly $C$ --- are detected rather than
resolved by luck. Floating point is used freely to \emph{propose} levels and
bisection points; it decides nothing.

The cost is set by the degree. With $d = d_g$, $\deg A = 2d$ and
$\deg B = d$, so $A'B$ and $2B'A$ both have degree $3d-1$ with leading
coefficient $2d\,\alpha_{2d}\beta_d$: the leading terms cancel, $\deg N \leq
3d-2$, and
\[
  \deg H \;=\; \deg B + \deg N \;\leq\; 4d - 2 .
\]
At $d = 11$ that is a chain on a degree-42 integer polynomial, which is why
the implementation carries the chain in GMP rather than in machine words.

Unwinding $H = BN$ gives the mechanisms by which the Morse property can fail.
A root of $H$ is repeated exactly when

\begin{enumerate}
\item[(M1)] $B$ has a repeated real root; or
\item[(M2)] $N$ has a repeated real root; or
\item[(M3)] $B$ and $N$ have a common real root.
\end{enumerate}

Since a root of $B$ is a point of the axis $a = 0$, (M1) and (M3) are
degeneracies located on that axis. Two further, non-generic, failures are not
of this form at all:

\begin{enumerate}
\item[(M4)] $B \equiv 0$, i.e.\ $f \perp g(b\,\cdot)$ for every $b$; then
  $u \equiv C$ and the whole backbone is critical.
\item[(M5)] $N \equiv 0$, i.e.\ $A'B = 2B'A$ identically, which integrates to
  $B^2 = cA$ for a constant $c$; then $u \equiv C - c$ and again the whole
  backbone is critical.
\end{enumerate}

The familiar example $f = g = x^m$ falls under (M5): the critical set is a
curve, not a finite set, and the Morse theory below does not apply to it.

\section{The Morse skeleton, and its one gap}\label{sec:skeleton}

Suppose $L$ is Morse. The exact enumeration of \S3 yields the finite critical
set together with the type of each point, and the critical values are
$u$ evaluated there. Because $L = u + Aw^2$ with $A > 0$, each sublevel set
$\{L \leq c\}$ is symmetric about the backbone and fibred over
$\{b : u(b) \leq c\}$ with interval fibres; its connected components are
therefore in bijection with the components of the one-dimensional sublevel set
$\{u \leq c\}$. The merge tree of $L$ --- which components are born at which
minima and merge at which saddles --- is thus determined combinatorially by
the ordering of the critical values of $u$, and is again exactly computable.

What the merge tree does not determine is the \emph{connectivity of the
invariant manifolds} of the gradient flow $\dot z = -\grad L(z)$. Each saddle
has a one-dimensional unstable manifold whose two branches descend, and a
one-dimensional stable manifold whose two branches ascend; the stable
manifolds partition the plane into the basins of the minima. Generically each
descending branch terminates at a minimum, and the resulting decomposition ---
the Morse--Smale skeleton --- is the object of interest. The exception is a
\emph{saddle connection}: a trajectory descending from one saddle directly into
another. When one is present the system is Morse but not Morse--Smale, the
stable and unstable manifolds coincide along an orbit, and the assignment of
branches to minima is not what the picture would suggest.

Saddle connections are codimension one in parameter space, so a given model
generically has none. But ``generically none'' is not a certificate, and
deciding the question for a \emph{particular} model is the remaining gap. The
next section shows that the gap cannot be closed by the algebraic methods of
\S\S2--3.

\section{The rheostat: an algebraically invisible bifurcation}\label{sec:rheostat}

Fix $f$, $g$ and $\mu$ and consider the one-parameter family obtained by
rescaling the two polynomials against each other,
\begin{equation}\label{eq:rheostat}
  f_\Lambda = \frac{f}{\sqrt\Lambda}, \qquad
  g_\Lambda = \sqrt\Lambda\, g, \qquad \Lambda > 0,
\end{equation}
with $\mu$ unchanged. We call $\Lambda$ the \emph{rheostat}. Write
$L_\Lambda$, $A_\Lambda$, $u_\Lambda$ for the objects built from
$(f_\Lambda, g_\Lambda, \mu)$, so that $L_1 = L$.

\begin{lemma}[Similarity]\label{lem:similarity}
For every $\Lambda > 0$ and every $(a,b) \in \R^2$,
\begin{equation}\label{eq:similarity}
  L_\Lambda(a, b) \;=\; \frac{1}{\Lambda}\, L\bigl(\Lambda a,\; b\bigr).
\end{equation}
\end{lemma}

\begin{proof}
From \eqref{eq:ABC}, $A_\Lambda = \Lambda A$, $B_\Lambda = B$ (the factors
$\sqrt\Lambda$ and $1/\sqrt\Lambda$ cancel) and $C_\Lambda = C/\Lambda$. Hence
\[
  L_\Lambda(a,b) = \frac{C}{\Lambda} - 2aB + a^2 \Lambda A
  = \frac{1}{\Lambda}\Bigl( C - 2(\Lambda a)B + (\Lambda a)^2 A \Bigr)
  = \frac{1}{\Lambda} L(\Lambda a, b). \qedhere
\]
\end{proof}

\begin{corollary}\label{cor:invariants}
Along the rheostat:
\begin{enumerate}
\item[(i)] the critical $b$-values are independent of $\Lambda$;
\item[(ii)] each critical point moves to $(a^*/\Lambda,\, b)$, i.e.\ slides
  along its vertical line towards the axis;
\item[(iii)] every critical value scales by the same positive factor,
  $u_\Lambda = u/\Lambda$, so the \emph{ordering} of the critical values, and
  hence the merge tree, is constant;
\item[(iv)] the type of each critical point is constant.
\end{enumerate}
\end{corollary}

\begin{proof}
$B_\Lambda = B$ and $N_\Lambda = A_\Lambda'B_\Lambda - 2B_\Lambda'A_\Lambda =
\Lambda N$, so $B_\Lambda N_\Lambda = \Lambda BN$ has the same roots as $BN$,
giving (i); $a^*_\Lambda = B/(\Lambda A) = a^*/\Lambda$ gives (ii); and
$u_\Lambda = C/\Lambda - B^2/(\Lambda A) = u/\Lambda$ gives (iii), the
ordering being preserved because $1/\Lambda > 0$. Finally $u_\Lambda'' =
u''/\Lambda$ has the sign of $u''$, which with
Proposition~\ref{prop:critical}(iii) gives (iv).
\end{proof}

Thus \eqref{eq:similarity} exhibits the family as an anisotropic dilation of
the plane, $(a,b) \mapsto (\Lambda a, b)$, composed with a rescaling of the
value. That map is a diffeomorphism, so it carries sublevel sets to sublevel
sets and preserves all of their topology: the merge tree, the number and type
of the critical points, the component structure, and the order of the critical
values are all constant along the rheostat.

The gradient flow, however, is not conjugated by it, and the reason is worth
stating precisely, because it identifies what the family actually is.

\begin{proposition}[The rheostat flow is a gradient-descent-like flow for a
fixed $L$]\label{prop:rsgd}
Let $T = \operatorname{diag}(\Lambda, 1)$ and let $z(t)$ solve $\dot z =
-\grad L_\Lambda(z)$. Then $\zeta = Tz$ satisfies
\begin{equation}\label{eq:preconditioned}
  \dot\zeta \;=\; -M\,\grad L(\zeta), \qquad
  M \;=\; \operatorname{diag}\bigl(\Lambda,\; \Lambda^{-1}\bigr),
\end{equation}
with $M$ symmetric positive definite and $\det M = 1$.
\end{proposition}

\begin{proof}
Differentiating \eqref{eq:similarity}, $\grad L_\Lambda(z) = \Lambda^{-1}
T\,\grad L(Tz) = \operatorname{diag}(1, \Lambda^{-1})\,\grad L(Tz)$, so
$\dot\zeta = T\dot z = -T\operatorname{diag}(1,\Lambda^{-1})\,\grad L(\zeta)
= -\operatorname{diag}(\Lambda, \Lambda^{-1})\,\grad L(\zeta)$.
\end{proof}

So the rheostat is not a family of unrelated landscapes: after the coordinate
change it is \emph{one} function $L$, flowed in a one-parameter family of
metrics. Since $M$ is positive definite, $\grad L^{\mathsf T}M\grad L > 0$ off
the critical set, so \eqref{eq:preconditioned} is a gradient-descent-like flow
for $L$; by the representation theorem for such flows it is in particular a
rotated--scaled gradient-descent flow $-\kappa^{-1}Q\,\grad L$, with $Q$ the
rotation carrying $\grad L$ to $-M\grad L$ and
$\kappa = \|\grad L\|/\|M\grad L\|$. Rotated--scaled descent is strictly more
general than preconditioning --- $Q$ is orthogonal with positive definite
symmetric part, so it carries a skew component moving \emph{along} the level
set, which no symmetric preconditioner produces --- and the rheostat lands in
the symmetric sub-case.

This makes Corollary~\ref{cor:invariants} an instance of a general fact rather
than an accident of the particular scaling. A gradient-descent-like flow has
the same critical points as the gradient flow, and at a critical point the
linearisation of \eqref{eq:preconditioned} is $-M\operatorname{Hess}L$, similar
to $-M^{1/2}(\operatorname{Hess}L)M^{1/2}$ and therefore, by Sylvester's law of
inertia, of the same signature as $-\operatorname{Hess}L$. Minima stay minima,
saddles stay saddles, and the merge tree is untouched --- because a change of
metric cannot move any of it.

What a change of metric does \emph{not} preserve is the trajectories. Two
gradient-descent-like flows for the same $L$ share the critical points and the
Morse data and have different orbits; and a saddle connection is a property of
an orbit, not of the critical data. That is the whole of the matter: the Morse
skeleton is metric-independent, the Morse--Smale skeleton is not.

\begin{remark}[The rheostat as a normalisation]
Since $L_1$ is the model as given, any model may be regarded as the point
$\Lambda = 1$ of its own rheostat, and the distance to the nearest saddle
connection along the family is $|\Lambda^* - 1|$. Because $f_\Lambda \approx
f\,(1 - (\Lambda-1)/2)$ near $\Lambda = 1$, that is a relative coefficient
perturbation of $|\Lambda^*-1|/2$, which makes the quantity directly comparable
with the precision to which the coefficients are known. It is of course a
distance along one direction only, hence an upper bound for the distance to the
non-Morse--Smale set; \S\ref{sec:certify} obtains a lower bound valid in all
directions.
\end{remark}

\subsection*{Why elimination cannot decide connectivity}

Corollary~\ref{cor:invariants} is the obstruction. Along the rheostat the
critical set $\{BN = 0\}$ is constant, so in $(b,\Lambda)$-space it is a union
of vertical lines: there is no discriminant to compute, no branching, nothing
for elimination theory to see. Yet the connectivity bifurcates. It follows that
no function of the coefficients of $A$, $B$ and $N$ --- no resultant, no
discriminant, no subresultant --- can decide whether a saddle connection is
present, since any such function is constant on a family along which the answer
changes.

This is not a statement that algebraic geometry is silent about the problem.
For a general one-parameter family, where the coefficients genuinely vary, the
set of parameters at which the fibre acquires a repeated critical point is the
discriminant locus of the curve $\{(b,\lambda) : (BN)(b;\lambda) = 0\}$, and it
is computed by elimination in one step rather than one parameter at a time.
That is exactly the \emph{non-Morse} set of \S3. The point is that the
\emph{non-Smale} set is a different object, invisible to the same machinery.

The correct setting is instead the one in which the splitting of invariant
manifolds is measured by an integral along the unperturbed connection. For a
regular level $c$ between the two critical values, let $\sigma(\Lambda)$ denote
the signed separation, measured along the level curve $\{L_\Lambda = c\}$,
between the point where the descending branch of the source saddle crosses it
and the point where the relevant ascending branch of the target saddle does.
Then $\sigma$ is analytic in $\Lambda$ wherever both crossings exist and are
transversal, and $\sigma(\Lambda) = 0$ is precisely a saddle connection. This
is a Melnikov function in all but name, and the expectation --- standard, and
the reason Hilbert's sixteenth problem is hard --- is that its zeros are
transcendental in the parameters. We are not aware of a proof of that
expectation for the present family and state it as an expectation only. What
is proved here is the weaker and sufficient statement of the preceding
paragraph: for the rheostat, no algebraic invariant can decide.

\section{Certifying the connection: existence}\label{sec:certify}

The rheostat is one path through a family of metrics on the plane, and
nothing in the existence argument below uses more than that. We state it for
a general such path, because the hypotheses separate cleanly: confinement is
a property of $L$ alone and is the same for every member of the family, while
only the endpoint landings depend on the particular path.

Let $\{M_s\}_{s \in [0,1]}$ be symmetric positive definite matrices depending
smoothly on $s$, and consider the flows
\begin{equation}\label{eq:family}
  \dot z \;=\; -M_s\,\grad L(z).
\end{equation}
By Proposition~\ref{prop:rsgd} the rheostat restricted to
$[\Lambda_0, \Lambda_1]$ is such a family, with $M_\Lambda =
\operatorname{diag}(\Lambda, \Lambda^{-1})$. Each member is
gradient-descent-like for $L$ (Appendix~\ref{app:rsgd}), which is all that is
used below.

\subsection*{Confinement does not depend on the metric}

The ingredient that makes the argument finite is that the orbit cannot escape,
and it holds for every descent flow at once.

\begin{lemma}[Vertical barrier]\label{lem:barrier}
For every $b_0$ and every $a$, $L(a, b_0) \geq u(b_0)$. Consequently an orbit
of any flow along which $L$ is non-increasing, and which is at level below
$u(b_0)$ at some time, can never cross the line $b = b_0$ thereafter.
\end{lemma}

\begin{proof}
$L(a,b_0) = u(b_0) + A(b_0)(a - a^*(b_0))^2 \geq u(b_0)$ since $A > 0$.
\end{proof}

In particular let $b_0$ be a simple real root of $B$. By the roof computation
of \S2, $u(b_0) = C$ and $b_0$ carries a saddle at that level. A branch
leaving a saddle of level $C$ has $L < C$ strictly for positive time, so by
Lemma~\ref{lem:barrier} it is trapped between the two nearest $B$-roots on
either side; and on that compact $b$-interval $A$ is bounded below, so
$L = u + A w^2$ bounded above bounds $w$, hence $a$, as well. The forward
orbit is contained in a compact set.

Nothing in this used the metric: only that $L$ decreases along the flow, which
holds for every member of \eqref{eq:family} since $M_s \succ 0$. The
confinement is a statement about the sublevel geometry of $L$, and the whole
family inherits it.

\subsection*{The limit is a single critical point}

A bounded orbit of \eqref{eq:family} converges to a single critical point of
$L$: this is Lemma~\ref{lem:kl}, the {\L}ojasiewicz inequality being a
property of $L$ that the flow inherits with its exponent and a constant
degraded only by $\lambda_{\min}(M_s)$.

So for each $s$ the branch in question has a single critical point as its
$\omega$-limit, and by strict decrease of $L$ that point lies strictly below
the source level and inside the confining interval. The candidates are a
finite, exactly enumerated list: the minima in the interval, and the saddles
in it of strictly lower level.

\subsection*{Intermediate values on the family}

\begin{theorem}[Existence of a connection]\label{thm:existence}
Let $\{M_s\}$ be as above, let the branch be confined as in
Lemma~\ref{lem:barrier} for every $s$, and suppose that at $s = 0$ and $s = 1$
the branch converges to two \emph{different} minima $m_0 \neq m_1$. Then there
is $s^* \in (0,1)$ at which the branch converges to a saddle: the flow
\eqref{eq:family} at $s^*$ has a saddle--saddle connection.
\end{theorem}

\begin{proof}
For a hyperbolic sink $m$, the set of $s$ for which the branch converges to
$m$ is open, by hyperbolicity together with smooth dependence of the unstable
branch on $s$. The sets $\{s : \text{limit} = m\}$, over the finitely many
minima in the confining interval, are therefore disjoint open sets. They
cannot cover $[0,1]$: $[0,1]$ is connected and $s=0$, $s=1$ lie in different
members of the family. Hence some $s^*$ has a limit that is not a minimum; by
the previous subsection the limit is a critical point, so it is a saddle.
\end{proof}

This is a topological argument. It uses no transversality, no continuity of a
shooting residual, no tracking of a section crossing across the interval, and
no control of the flow at infinity --- the barrier removes the far field
before compactification is needed. (SPONG's Poincar\'e compactification is
available and its equatorial equilibria are known, but the degenerate
backbone ends carry only asymptotic-grade certificates, so an argument that
avoids them is worth preferring.)

What the theorem does \emph{not} give is simplicity or uniqueness of the wall.
It produces at least one $s^*$; the family may cross several, and the
connection at a given $s^*$ need not be the only one. Any claim that a
particular measured crossing is a single simple handle slide is a further
hypothesis, requiring transversality and an exclusion argument for
flip-and-return pairs between sample points.

The two hypotheses consumed are the endpoint landings. Each is a capture deep
inside a sink basin --- robust, but a numerical observation. The next section
is what upgrades them.

\section{Certifying the absence of a connection}

[To be written: separation of the two manifolds at a common regular level by
Gr\"onwall enclosure of the propagated launch data; the resulting ball in
coefficient space within which no connection exists; measured margins; and
the backward-stability licence under which a connection may be drawn.]

\section{Rendering the invariant manifolds}

[To be written: why na\"\i ve integration fails here, anadromic schemes, the
Hadamard graph transform, backbone guidance, and the comparison against
forward/backward Euler, RKF45 and Stork.]

\section{Descent methods on a known landscape}

[To be written: the mean-field ODE is the portrait itself; Adam follows a
different vector field (Dereich--Jentzen--Kassing) which nevertheless shares
the orbits in the noiseless limit; the non-interpolating regime; basins and
the case for many initialisations.]

\section{Allocating work across initialisations}

[To be written: the bandit formulation and what it buys when the basins are
unknown.]

\appendix

\section{Descent flows in a fixed metric}\label{app:rsgd}

This appendix proves the two facts about descent flows used in
\S\ref{sec:rheostat} and \S\ref{sec:certify}. Both are elementary for the
flows the note actually needs, in which the multiplier is a constant
symmetric positive definite matrix; Remark~\ref{rem:rsgd} says what changes
in general.

\begin{definition}[Gradient-descent-like]\label{def:gdl}
Let $L \in C^2(\R^d)$ and $\Omega = \{y : \grad L(y) \neq 0\}$. A $C^1$ field
$v$ is \emph{gradient-descent-like} for $L$ if $v^{\mathsf T}\grad L < 0$ on
$\Omega$ and $v$ vanishes exactly on the critical set of $L$.
\end{definition}

The first condition says $L$ strictly decreases along the flow: $v$ need not
be steepest descent, only descent. The second aligns the equilibria of $v$
with the critical points of $L$. For $v = -M\grad L$ with $M$ constant and
symmetric positive definite --- the case of Proposition~\ref{prop:rsgd},
where $M = \operatorname{diag}(\Lambda, \Lambda^{-1})$ --- both hold at once,
since $\grad L^{\mathsf T}M\grad L > 0$ off the critical set and $M$ is
invertible.

\begin{lemma}[Inertia]\label{lem:inertia}
Let $M$ be symmetric positive definite and $H$ symmetric. Then $MH$ has the
same numbers of positive, negative and zero eigenvalues as $H$.
\end{lemma}

\begin{proof}
$MH$ is similar to $M^{-1/2}(MH)M^{1/2} = M^{1/2}HM^{1/2}$, which is symmetric
and congruent to $H$; apply Sylvester's law of inertia.
\end{proof}

With $H = -\operatorname{Hess}L(y_*)$ this says the flow $\dot y = -M\grad L$
has the same numbers of attracting and repelling directions at every critical
point as the gradient flow. When $L$ is Morse the Morse indices agree: minima
remain sinks, saddles remain saddles. This is the step used in
\S\ref{sec:rheostat}.

\begin{lemma}[{\L}ojasiewicz]\label{lem:kl}
Let $L$ be real-analytic and let $y(t)$ solve $\dot y = -M\grad L(y)$ with $M$
constant symmetric positive definite. Then along the flow
\begin{equation}\label{eq:klrate}
  -\frac{d}{dt}L(y(t)) \;=\; \grad L^{\mathsf T}M\,\grad L
  \;\geq\; \lambda_{\min}(M)\,\|\grad L(y(t))\|^2,
\end{equation}
and every bounded orbit converges to a single critical point of $L$.
\end{lemma}

\begin{proof}
\eqref{eq:klrate} is immediate. Near a critical point $y_*$, analyticity gives
a {\L}ojasiewicz inequality $|L - L(y_*)|^{\beta} \leq C\|\grad L\|$ for some
$\beta \in [\tfrac12, 1)$, so \eqref{eq:klrate} yields
$-\tfrac{d}{dt}L \geq \lambda_{\min}(M)C^{-2}|L - L(y_*)|^{2\beta}$: the
inequality is a property of $L$, inherited with its exponent and a constant
degraded only by $\lambda_{\min}(M)$. The usual length argument then applies
--- $\|\dot y\| \leq \lambda_{\max}(M)\|\grad L\|$ bounds the arclength by a
convergent integral --- so the orbit has finite length and its $\omega$-limit
is a single point, which is critical by Definition~\ref{def:gdl}.
\end{proof}

\begin{remark}\label{rem:rsgd}
Both lemmas hold more generally for a variable multiplier $M(y)$, symmetric
positive definite and $C^1$: the linearisation is then no longer of the form
$MH$ with $M$ constant, so the inertia statement requires the
Ostrowski--Schneider theorem rather than Sylvester's law, and the
{\L}ojasiewicz constant degrades by $\lambda_{\min}(M(y_*))$ locally rather
than globally. Neither generalisation is needed here, where $M$ is constant
and diagonal.

The parameterisation $v = -\kappa^{-1}Q\,\grad L$ with $\kappa > 0$ and $Q$
orthogonal of positive definite symmetric part --- rotated--scaled gradient
descent --- realises the same descent fields with a multiplier of condition
number one, which matters in floating point and not on the blackboard; see
the author's separate work on those flows.
\end{remark}

\end{document}
